\documentclass{beamer}

\usepackage[utf8]{inputenc}
\usepackage[french]{babel}

\usepackage{verbatim}
\usepackage{minted}

\usepackage{tikz}

\newcommand{\tocpage}{
  \begin{frame}
    \tableofcontents[currentsection]
  \end{frame}
}

\def\labelenumi{\theenumi}
\usepackage{enumitem}

\setitemize[1]{label=$\circ$}
%\usetheme{Luebeck}
\usetheme{CambridgeUS}

% \usecolortheme{dolphin}
\usecolortheme{seahorse}

\setbeamertemplate{section in toc}{%
   ({\inserttocsectionnumber})~{\inserttocsection}%
}

\setbeamertemplate{navigation symbols}{}



\title{Présentation projet Fouine}
\author{Robin Jourde \& Nicolas Nardino}
\date{mai 2021}

\begin{document}


\begin{frame}
  \titlepage
\end{frame}


\begin{frame}{Dites \og Bonjour \fg~à notre amie la fouine !}
  \begin{center}
    \tiny
    \verbatiminput{fouine-ascii.txt}
  \end{center}
\end{frame}


\begin{frame}{Table de matières}
  \tableofcontents
\end{frame}


\section{Fouine, qu'est-ce que c'est ?}

\tocpage

\begin{frame}{Un terme polysémique}
  \begin{itemize}
  \item<2-> un petit animal trop mignon : \\
    \includegraphics[width=0.2\linewidth]{figures/photo_fouine.jpg}
  \item<3-> un sous langage de OCaml
  \item<4-> un interpréteur de ce sous langage
  \end{itemize}

  \visible<5-> {Par la suite on n'utilisera le terme que dans ses 2 derniers sens.}
\end{frame}

\begin{frame}{Le langage}
  Grammaire simplifiée :
  \begin{align*}
    e~:=&~let~(rec)~p~=~e_1~in~e_2~|~fun~p~\to~e_1~|~e_1~e_2~|~e_1~:=~e_2 \\
        &|~!e_1~|~if~e_1~then~e_2~else~e_3~|~match~e_1~with~[|~p_i~->~e_i] \\
        &|~e_1~::~e_2~|~try~e_1~with~E~p~->~e_2~|~raise~e_1~|~p
  \end{align*}

  $p$ représente un pattern :
  $$
  p~:=~x~|~\_~|~k~|~p_1,~p_2~|~p_1~::~p_2
  $$

  $k$ est une constante

  \visible <2-> {Et les opérations arithmétiques me diriez vous ? ... }

  \visible <3-> {... Et là vous croyiez que j'allais le dire maintenant mais en fait non, il a falloir attendre ! Suspens ...!!  }
\end{frame}


\begin{frame}{L'interprète}
  \begin{itemize}

  \item<1-> \texttt{./fouine [options] test1515.ml}

  \item<2->  Mais quelles options ??? Demandons à fouine : \texttt{./fouine --help}

  \item<3->  \texttt{
      -showsrc -debug -tree -outval -stdin -cps -outcode -outcode-tree -run -autotest -optim -reduc -prefix -notypes -showtypes -showmoretypes -monotypes -cpstypes
    }

  \end{itemize}
\end{frame}


\begin{frame}{Un exemple}
  \begin{itemize}
  \item fichier de test \texttt{not.ml}

  \mint{ocaml}|
  let x = not (0>1) in if x then prInt 1 else prInt 0
  |

 \item
\texttt{ ./fouine -debug -outval -showtypes not.ml }

\texttt{ x : bool }

\texttt{ - : int }

\texttt{ let x = not (( > ) (0) (1)) in (if x then prInt (1) else prInt (0)) }

\texttt{ 1 }
 
\texttt{ 1 }


\end{itemize}
\end{frame}

\section{Fouine, comment ça marche ?}

\tocpage

\begin{frame}{Que fait fouine ?}
  Étapes :
  \begin{enumerate}
  \item Lire : lexer et parser
  \item Typer
  \item (Si souhaité) traduction CPS
  \item (Si souhaité) réduction
  \item Évaluer
  \end{enumerate}

  Pour tout cela on a besoin de représenter les expressions...
\end{frame}

\begin{frame}{Des constructeurs et des types}
  \usetikzlibrary {arrows.meta,graphs,graphdrawing}
\usegdlibrary {layered}

\tiny
\tikz[>={Stealth[round,sep]}]
\graph [layered layout, sibling distance = 1pt, layer distance = 0.8cm,
nodes=draw, edges=rounded corners]
{

  const -> {
    Int ,
    Unit ,
    Bool ,
    Nil
  } ;

  value -> {
    Const -> const  ,
    Ref ,
    StdLib -> expr  ,
    VFun -> { pattern , env , expr } ,
    VCpl  ,
    VList
  } ;
  value <- {
    StdLib ,
    VCpl ,
    VList
  } ;
  env -> value ;

  pattern -> {
    Ident ,
    Under ,
    PConst -> const ,
    PCpl ,
    PList
  } ;
  pattern <- {,
    PCpl ,
    PList
  } ;

  expr -> {
    Val -> value ,
    Pattern -> pattern ,
    Let -> pattern ,
    Fun -> pattern ,
    Rec -> pattern,
    App ,
    Cpl ,
    Aff ,
    Der ,
    If ,
    Match -> pattern matching ,
    Cons ,
    Try -> pattern ,
    Raise
  };
  expr <- {
    Let ,
    Fun ,
    Rec ,
    App ,
    Cpl ,
    Aff ,
    Der ,
    If ,
    Match ,
    Cons ,
    Try ,
    Raise
  };

  pattern matching -> { pattern, expr } ;

};

\end{frame}



\begin{frame}{La bibliothèque standard}
  On tient notre promesse !

  \begin{itemize}
  \item \texttt{fst , snd}
  \item \texttt{prInt}
  \item \texttt{ref}
  \item \texttt{( + ) , ( - ) , ( * ) , ( / )}
  \item \texttt{( \&\& ) ,  ( || )}
  \item \texttt{not}
  \item \texttt{( <= ) , ( >= ) , ( < ) , ( > )}
  \item \texttt{( = ) , ( <> )}
  \end{itemize}
\end{frame}

\begin{frame}{Exemple détaillé : \texttt{prInt}}
  \begin{itemize}
  \item code de la fonction

    \texttt{ let \_prInt m = function }

    \texttt{ | Const (Int i) -> (print\_int i ; print\_newline () ; \~\& (Int i) ) }

    \texttt{ | \_ -> raise PrInt\_not\_int }

  \item insertion dans l'environnement

    \mint{ocaml}|("prInt", StdLib _prInt) :: ...|

  \item type

    \mint{ocaml}|("prInt", TFun(TInt, TInt) ):: ... |


  \end{itemize}
  \end{frame}

\section{Conclusion}
\begin{frame}{C'est fini !}
  \begin{center}
    \Huge Merci à tous \& Bisous !
  \end{center}
\end{frame}

\begin{frame}
  \begin{center}
    \fontsize{2}{2} \selectfont
    \verbatiminput{daniel-ascii.tex}
  \end{center}
\end{frame}

\end{document}



